\begin{lemma}[Original section coefficient naturality]
\label{scheme:actual-section-coefficient-naturality}
\uses{scheme:relative-section-coefficients}
\lean{EnvelopeProperty.SchemeImage.sectionBaseMap_comp_app}
\leanok
Let $R$ and $K$ be any commutative rings, let $\varphi:R\to K$ be an
actual ring homomorphism, and let $g:X\to Y$, $f_X:X\to\operatorname{Spec}K$,
$f_Y:Y\to\operatorname{Spec}R$ be original scheme morphisms satisfying
$g f_Y=f_X\operatorname{Spec}(\varphi)$. For every actual open $U\subset Y$,
the whole original coefficient maps satisfy
\[
 g^*\circ\operatorname{sectionBaseMap}(f_Y,U)
 =\operatorname{sectionBaseMap}(f_X,g^{-1}U)\circ\varphi.
\]
There is no field, injectivity, flatness, finite-type or reducedness
assumption.
\end{lemma}
\begin{proof}
Fix $r\in R$ and regard it as the original global section on
$\operatorname{Spec}R$ under the canonical affine global-section
isomorphism. Pullback of that section along the actual commuting scheme
triangle gives
$g^*(f_Y^*(r))=f_X^*(\varphi(r))$, because the affine global-section
isomorphism is natural with respect to the actual homomorphism $\varphi$.
Restrict the equality to $g^{-1}U$. Naturality of the original sheaf
pullback identifies restriction after the global pullback along $g$ with
pullback along $g$ of the restriction to $U$. These are precisely the
maps in the definition of the two named section coefficient maps. Thus
the displayed equality holds on every $r$, hence as a whole ring-map
identity. This is the existing generic naturality theorem, isolated from
its later application to flat scheme images.
\end{proof}
